% ----- CHAUPAI 18 -----

\newpage

\subsection*{\deva{अष्टादश चौपाई} \(Eighteenth Chaupai\)}
\addcontentsline{toc}{subsection}{\deva{अष्टादश चौपाई} \(Eighteenth Chaupai\) — \deva{जुग सहस्र...}}

\begin{minipage}[t]{0.55\textwidth}
\vspace{0pt}

{\large \linespread{1.5}\selectfont
\deva{जुग सहस्र जोजन पर भानू।}\\
\deva{लील्यो ताहि मधुर फल जानू॥}
\par}

\vspace{0.5em}

{\large \linespread{1.5}\selectfont
\telu{జుగ సహస్ర జోజన పర భానూ।}\\
\telu{లీల్యో తాహి మధుర ఫల జానూ॥}
\par}

\vspace{0.5em}

{\large \linespread{1.3}\selectfont
\textit{juga sahasra jojana para bhānū |}\\
\textit{līlyo tāhi madhura phala jānū ||}
\par}
\end{minipage}%
\hfill
\begin{minipage}[t]{0.42\textwidth}
\vspace{0pt}
\begin{tcolorbox}[anchorbox]
\textbf{Fearless Curiosity:} As a child, he literally leaped for the sun because he wanted to understand it. Fearless curiosity and ambition are the primary engines of learning. Students should never be afraid to reach high or ask bold questions, even if they occasionally make mistakes along the way.
\vspace{0.5em}\newline Driven by fearless childhood curiosity, he leaped into the unknown, teaching us to never fear making bold mistakes while learning.\newline\null\hfill\href{https://www.valmiki.iitk.ac.in/sloka?field_kanda_tid=4&language=dv&field_sarga_value=66}{\textit{Kishkindha Kanda, 66:21-23}}
\end{tcolorbox}
\end{minipage}

\vspace{0.5em}
\subsubsection*{\textbf{Visual Rhythm Map:}}
\begin{table}[h!]
\centering
\renewcommand{\arraystretch}{1.2}
\setlength{\tabcolsep}{0pt}
\begin{tabularx}{\textwidth}{|*{16}{>{\centering\arraybackslash\hsize=1.0\hsize}X|}}
\hline
\multicolumn{2}{|>{\centering\arraybackslash\hsize=2.0\hsize}X|}{\textbf{\laghu{जु}\laghu{ग}} \par (2)} &
\multicolumn{4}{>{\centering\arraybackslash\hsize=4.0\hsize}X|}{\textbf{\laghu{स}\guru{ह}\laghu{स्र}} \par (4)} &
\multicolumn{4}{>{\centering\arraybackslash\hsize=4.0\hsize}X|}{\textbf{\guru{जो}\laghu{ज}\laghu{न}} \par (4)} &
\multicolumn{2}{>{\centering\arraybackslash\hsize=2.0\hsize}X|}{\textbf{\laghu{प}\laghu{र}} \par (2)} &
\multicolumn{4}{>{\centering\arraybackslash\hsize=4.0\hsize}X|}{\textbf{\guru{भा}\guru{नू}} \par (4)} \\
\hline
\end{tabularx}
\vspace{-1.5em}
\end{table}

\begin{table}[h!]
\centering
\renewcommand{\arraystretch}{1.2}
\setlength{\tabcolsep}{0pt}
\begin{tabularx}{\textwidth}{|*{16}{>{\centering\arraybackslash\hsize=1.0\hsize}X|}}
\hline
\multicolumn{4}{|>{\centering\arraybackslash\hsize=4.0\hsize}X|}{\textbf{\guru{ली}\guru{ल्यो}} \par (4)} &
\multicolumn{3}{>{\centering\arraybackslash\hsize=3.0\hsize}X|}{\textbf{\guru{ता}\laghu{हि}} \par (3)} &
\multicolumn{3}{>{\centering\arraybackslash\hsize=3.0\hsize}X|}{\textbf{\laghu{म}\laghu{धु}\laghu{र}} \par (3)} &
\multicolumn{2}{>{\centering\arraybackslash\hsize=2.0\hsize}X|}{\textbf{\laghu{फ}\laghu{ल}} \par (2)} &
\multicolumn{4}{>{\centering\arraybackslash\hsize=4.0\hsize}X|}{\textbf{\guru{जा}\guru{नू}} \par (4)} \\
\hline
\end{tabularx}
\vspace{-1.5em}
\end{table}

\vspace{0.5em}
\textbf{ \deva{सरल-संस्कृतम्}:}\\
 \deva{युग-सहस्र-योजन-दूरे स्थितं भानुं (सूर्यम्)}\\
 \deva{मधुरं फलम् इति मत्वा भवान् लीलया निगीर्णवान् (भक्षितवान्)।}\\

The sun, situated thousands of Yojanas away \\
thinking it to be a sweet fruit, you playfully swallowed it.


\vspace{1em}
\newpage
\textbf{\deva{प्रतिपदार्थः} (Meaning of each word):}
\begin{table}[h!]
\centering
\renewcommand{\arraystretch}{1.2}
\begin{tabularx}{\textwidth}{@{} l l X X @{}}
\toprule
\textbf{अवधी} & \textbf{संस्कृतम्} & \textbf{English} & \textbf{Grammar \& Notes} \\
\midrule
\deva{जुग} \glossaryicon / \telu{జుగ} / \textit{juga}  &  \deva{युग}  &  Yuga (Era/Distance)  & \\ 
\deva{सहस्र} / \telu{సహస్ర} / \textit{sahasra}  &  \deva{सहस्र}  &  Thousand  &   \\
\deva{जोजन} / \telu{జోజన} / \textit{jojana}  &  \deva{योजन}  &  Yojana (approx 8 miles)  & \\ 
\deva{पर} / \telu{పర} / \textit{para}  &  \deva{दूरे / परत्र}  &  Away / At a distance  &   \\
\deva{भानू} / \telu{భానూ} / \textit{bhānū}  &  \deva{भानुम्}  &  The Sun  & Long `ū` for rhyme \\
\deva{लील्यो} / \telu{లీల్యో} / \textit{līlyo}  &  \deva{लीलया निगीर्णवान्}  &  Swallowed playfully  &   \\
\deva{ताहि} / \telu{తాహి} / \textit{tāhi}  &  \deva{तम्}  &  It (the sun)  & Awadhi objective \\
\deva{मधुर फल} / \telu{మధుర ఫల} / \textit{madhura phala}  &  \deva{मधुर-फलम्}  &  Sweet fruit  &   \\
\deva{जानू} / \telu{జానూ} / \textit{jānū}  &  \deva{ज्ञात्वा / मत्वा}  &  Knowing / Thinking  & Long `ū` for rhyme \\
\bottomrule
\end{tabularx}
\end{table}

\begin{tcolorbox}[poetrybox]
\textbf{The Verse the Internet Loves — The 96-Million-Mile Question:}\\
This is the most-quoted verse of the Hanuman Chalisa in modern popular science discussions. The calculation goes:
\deva{जुग} (\textit{Yuga}) = 12,000 $\times$ \deva{सहस्र} = 1,000 $\times$ \deva{जोजन} $\approx$ 8 miles:
\[
12{,}000 \times 1{,}000 \times 8~\text{miles} ~=~ 96{,}000{,}000~\text{miles}
\]
The modern scientific distance from the Earth to the Sun is approximately 93,000,000 miles ($\approx$ 150 million km). This similarity has made the verse enormously popular.

A few points of honest scholarly context:
\begin{itemize}
  \item The Hanuman Chalisa is a devotional poem, not an astronomical treatise. Neither it nor Tulsidas's Ramcharitmanas makes an explicit claim of astronomical accuracy.
  \item \deva{जुग} most commonly refers to a unit of time, not a distance scalar.  
    It can also mean \textit{joining}, allowing us to interpret the verse as a praise of Hanuman for joining with the sun which is thousands of \textit{yojanas} away.
  \item We cannot determine whether Tulsidas intentionally encoded the distance between the Earth and the Sun, or whether this is a beautiful numerical coincidence.  
  The poet could have used common knowledge fo the time, or he might have used it to refere to a large number, while fitting the metrical requirements of the Chalisa. 
\end{itemize}
 Neither interpretation diminishes the devotional or literary power of the verse.
\end{tcolorbox}

\newpage
